\section{Formal model and implementation boundary}

\subsection{Structural states and mutation}
Let $\mathcal{Z}$ be the set of encoded candidate instances and let
$\phi:\mathcal{Z}\to\mathcal{X}$ map an encoding to the solver-visible
structural phenotype. Encodings $z,z'$ are equivalent when
$\phi(z)=\phi(z')$. A structural state is a phenotype $x\in\mathcal{X}$,
for example a canonical graph/instance together with whichever structural
attributes determine solver input. The exact canonicalization and attributes
are implementation questions, not assumptions about FOGA-11.

For finite-budget analysis, take $\mathcal{X}_B$ to be the set of
phenotypes reachable within budget $B$; assume it is finite. A mutation kernel
$K_t(x'\mid x,s)$ is the conditional distribution of the next phenotype given
parent phenotype $x$ and full pre-mutation state $s$. It includes the selected
structural edit, repair, and canonicalization. Dependence on $t$ or $s$ permits
operator adaptation and history-dependent mutation. The corresponding kernel
on encodings is not necessarily Markov after projecting through $\phi$.

\subsection{Population memory and parent survival}
At proposal epoch $t$, let $P_t$ be the parent population (a multiset of
encodings or phenotypes) and $A_t$ any quality-diversity archive. Let $S_t$
collect scores, lineage information, solver state, and other algorithm
variables that affect the next transition. Thus $M_t=(P_t,A_t,S_t)$ is the
operational memory state. Parent survival is a transition event: for parent
$i$, its survival probability is
\[
 r_{t,i}=\Prb(i\in P_{t+1}\mid M_t,\text{environment at }t).
\]
This is distinct from its parent-selection probability
$\pi_{t,i}=\Prb(i\text{ is selected for a proposal}\mid M_t)$.
Retaining a parent does not imply it is selected, and selection does not
guarantee a useful mutation. The full process may be time-inhomogeneous
Markovian in $(M_t,\theta_t,q_t)$ if that state is sufficient; a niche label
or niche-count vector alone need not be.

\subsection{Hardness and moving target}
Let $\theta_t$ describe the task/environment and $q_t$ the state of the
co-adapting BQ solver. Define true hardness $h_t(x)=h(x;\theta_t,q_t)$,
where a larger value means harder under a prespecified solver-performance
metric. For threshold $\kappa$, the true hard region at epoch $t$ is
\[
 \HARD_t=\{x\in\mathcal{X}_B:h_t(x)\geq\kappa\}.
\]
This definition allows the target to move because either the task or solver
state changes. The event of discovery is evaluated at a declared evaluation
epoch: a proposal $X_t$ hits when $X_t\in\HARD_t$. For the proofs, $\HARD_t$
must be determined before drawing/evaluating $X_t$ (it may be random but
predictable from the pre-proposal history). Whether actual BLADE labels a
candidate ``hard'' by HGB40 alone or by BQ verification is unspecified.

HGB40 supplies an estimate $\widehat h_t(x)$, not the definition of the
latent target. Surrogate ranking can be used by selection, but an estimate of
reachability into $\HARD_t$ requires calibration or a direct verification
model. Without such a condition, true-hardness hit probabilities cannot be
inferred from HGB40 scores.

\subsection{First hit and computational normalization}
Index proposal/evaluation cycles by $i$. Let
\[
 \tau=\inf\{i\geq 1:X_i\in\HARD_i\},\qquad
 C_\tau=\sum_{i=1}^{\tau}c_i,
\]
where $c_i>0$ is the normalized compute cost through evaluating proposal $i$.
The budget-$B$ discovery event is $\{C_\tau\leq B\}$. The cost must include
the work that differs between methods (structural generation, HGB40 inference,
archive operations, and BQ evaluation/verification, as applicable). Counting
generations alone is not computational normalization. The elementary bound
below uses deterministic proposal costs $c_i$; random or method-dependent
costs require a cost-aware hazard analysis and are not silently replaced by
generation counts.

For a policy $\pi$, define its conditional hit hazard
\[
 q_i^\pi=\Prb(X_i\in\HARD_i\mid\mathcal{F}_{i-1}^\pi),
\]
where $\mathcal{F}_{i-1}^\pi$ is the complete pre-proposal history. On a
no-hit history, a lower bound on $q_i^\pi$ is a reachability certificate. A
Markov state representation is one way to calculate this hazard, but the
history-based definition does not require Markovianity.

\subsection{Abstraction versus supplied BLADE-NCO context}
\begin{center}
\begin{tabular}{>{\raggedright\arraybackslash}p{0.22\linewidth}>{\raggedright\arraybackslash}p{0.35\linewidth}>{\raggedright\arraybackslash}p{0.34\linewidth}}
\toprule
Supplied component & Representation here & Unspecified for a BLADE-specific proof \\
\midrule
FOGA-11 mutation & Phenotype kernel $K_t$ & Edit set, operator probabilities, repairs, and transition support \\
Inherited parents; Elite-QD1000 (P=T) & Parent population $P_t$ and archive $A_t$ & Meaning of $P=T$, archive update, parent survival, and selection probabilities \\
HGB40 and BQ co-adaptation & Surrogate $\widehat h_t$ and true $h(x;\theta_t,q_t)$ & Hardness metric/threshold, calibration, solver update timing, and validation cost \\
Fixed generator and HAC & Comparator proposal kernels & Exact proposal laws, gradient access, stopping rules, and per-proposal compute \\
Finite budget & Cumulative cost $C_\tau$ & Which compute resources are charged and how variable costs are normalized \\
\bottomrule
\end{tabular}
\end{center}
Accordingly, the formal model is a deliberately explicit abstraction, not a
claim that BLADE-NCO has exactly these transition probabilities or state
variables. The missing implementation semantics are recorded as open
assumptions in \texttt{assumption\_register.md}.

\subsection{Niche aggregation and the Markov condition}
Suppose $g$ maps the full state to a niche label or niche-count summary. The
aggregated process $g(M_t)$ is Markov only if states with the same aggregate
have equal transition probabilities into every aggregate block (strong
lumpability). Otherwise the aggregate hides parent identity, within-niche
structure, or solver/archive history that predicts the next transition. An
exact failure example is given in Section~\ref{sec:counterexamples}.